Define the Krull dimension of a ring, heights of ideals and Krull's principal ideal theorem

We say that a chain of prime ideals of the form
\(\mathfrak{p}_0\subsetneq \mathfrak{p}_1\subsetneq \ldots \subsetneq \mathfrak{p}_n\)
has length \( n \). That is, the length is the number of strict inclusions, not the number of primes; 
these differ by \( 1 \). We define the Krull dimension of \( R \) to be the supremum of the lengths of all chains of prime ideals in \( R \).

Given a prime ideal \(\mathfrak{p}\) in \(R\), we define the height of \(\mathfrak{p}\), written \(\operatorname{ht}(\mathfrak{p})\), 
to be the supremum of the lengths of all chains of prime ideals contained in \(\mathfrak{p}\).
In other words, the height of \( \mathfrak{p} \) is the Krull dimension of the localization of \( R \) at \( \mathfrak{p} \).


A prime ideal has height zero if and only if it is a minimal prime ideal.
The Krull dimension of a ring is the supremum of the heights of all maximal ideals, or those of all prime ideals. 
The height is also sometimes called the codimension, rank, or altitude of a prime ideal.
More generally, the height of an ideal \(I\)  is the infimum of the heights of all prime ideals containing \(I\).  
In the language of algebraic geometry, this is the codimension of the subvariety of \(\mathrm{Spec}(R)\) corresponding to \(I\).

For an \( R \)-module \( M \) the Krull dimension is defined as 
\( \text{dim}_R M = \text{dim}(R / \text{Ann}_R(M) \)
this allows to give in the context of schemes, to give a dimension for coherent sheaves, or generalized finite rank vector bundles.

Some useful facts:

For an affine variety its dimension is defined as the maximum dimension of its irreducible components.
Similarly for a projective variety this is the maximal dimension of its affine irreducible components.

Let \( R \) be an algebra over a field \( k \) that is an integral domain. 
Then the Krull dimension of \( R \) is less than or equal to the transcendence degree of the field of fractions of \( R \) over \( k \). 
Equality holds if \( R \) is finitely generated as an algebra (for instance by the Noether normalization lemma).

Irreducible polynomials in a noetherian \( R \) induce principal prime ideals of height \( 1 \).
A local ring has Krull dimension \( 0 \) if and only if every element of its maximal ideal is nilpotent.

Every Artinian ring (e.g. fields) has dimension \( 0 \)
Further for Noetherian/valuation ring \( R \) one has
\( \text{dim } R[x] = \text{dim } R + 1 \)
further for any prime \( \mathfrak{p} \in \text{Spec } R \)
\[
    \text{ht}(\mathfrak{p}R[x]) = \text{ht} ( \mathfrak{p} )
    \text{ht}(\mathfrak{q}) = \text{ht} ( \mathfrak{p} ) + 1
\]
for any \( \mathfrak{q} \supseteq \mathfrak{p}R[x] \) in \( R[x] \) that contracts to \( \mathfrak{p} \),
moreover the injection \( R \hookrightarrow R[x] \) enforces that each fiber \( \text{Spec } R[x] \to \text{Spec } R \)
has no chain of prime ideals of length \( \geq 2 \).

Warning:
In a Noetherian ring, every prime ideal has finite height, however there are Noetherian rings of infinite dimension.

Krull's principal ideal theorem

Precisely, if \( R \) is a Noetherian ring and \( R \) is a principal, proper ideal of \( R \), then each minimal prime ideal containing \( I \) 
has height at most one.

This theorem can be generalized to ideals that are not principal, and the result is often called Krull's height theorem.  
This says that if \( R \) is a Noetherian ring and \( I \) is a proper ideal generated by \( n \) elements of \( R \), 
then each minimal prime over \( I \) has height at most \( n \). 
The converse is also true: if a prime ideal has height \( n \), then it is a minimal prime ideal over an ideal generated by \( n \) elements.

